% =====================================================================
%  outline.tex —— 文档结构总览
%  这一章是本项目"模块化"主张在文档侧的自述：9 个插件 -> 9 章，一一对应。
%  同时诚实标注当前成文进度。
% =====================================================================

\chapter*{文档结构总览}
\addcontentsline{toc}{chapter}{文档结构总览}
\markboth{文档结构总览}{文档结构总览}

本文档与代码仓库的 \fpath{huntian/plugins/} 目录一一对应：\textbf{一个插件一个模块章}。
模块章的编号、文件名与模块标识的对应关系是稳定的契约，后续版本只增不改。

\vspace{0.5em}
\small
\begin{longtable}{@{}c l l l c@{}}
\toprule
\textbf{章} & \textbf{模块标识} & \textbf{版本} & \textbf{类别} & \textbf{状态}\\
\midrule
\endfirsthead
\toprule
\textbf{章} & \textbf{模块标识} & \textbf{版本} & \textbf{类别} & \textbf{状态}\\
\midrule
\endhead
\bottomrule
\endfoot

\multicolumn{5}{@{}l@{}}{\itshape 第一部分\ \ 总论}\\
1  & 引言与文档范围                       & ---     & ---        & 待写\\
2  & 系统架构：模块即插件、服务即插件       & ---     & ---        & 待写\\
3  & 数据流与 artifact 契约               & ---     & ---        & 待写\\
4  & 验证方法学                           & ---     & ---        & 待写\\
5  & 模块索引与依赖关系                    & ---     & ---        & 待写\\
\midrule
\multicolumn{5}{@{}l@{}}{\itshape 第二部分\ \ 网格与材料}\\
6  & \pid{mesh.structured2d}             & 0.1.0   & mesh       & \textbf{已成文}\\
7  & \pid{materials.dft}                 & 0.1.0   & physics    & \textbf{已成文}\\
\midrule
\multicolumn{5}{@{}l@{}}{\itshape 第三部分\ \ 极化与量子}\\
8  & \pid{physics.polarization.ambacher} & 0.1.0   & physics    & \textbf{已成文}\\
9  & \pid{physics.quantum.analytic_tw}  & 0.1.0   & physics    & \textbf{已成文}\\
10 & \pid{physics.quantum.sp1d}          & 1.1.1   & physics    & \textbf{已成文}\\
\midrule
\multicolumn{5}{@{}l@{}}{\itshape 第四部分\ \ 输运与提取}\\
11 & \pid{transport.dd}                  & 0.1.0   & transport  & \textbf{已成文}\\
12 & \pid{extract.fom}                   & 0.1.0   & extract    & \textbf{已成文}\\
\midrule
\multicolumn{5}{@{}l@{}}{\itshape 第五部分\ \ 数值服务模块}\\
13 & \pid{nonlinear.newton}              & 1.0.0   & solver     & \textbf{已成文}\\
14 & \pid{linear.blocktridiag}           & 0.1.0   & solver     & \textbf{已成文}\\
\midrule
\multicolumn{5}{@{}l@{}}{\itshape 附录}\\
A  & 闭式解基准汇总                       & ---     & ---        & 待写\\
B  & 可标定参数总表                       & ---     & ---        & 待写\\
C  & 版本、指纹与 artifact 契约矩阵        & ---     & ---        & 待写\\
D  & 已知边界与欠账清单                    & ---     & ---        & 待写\\
\midrule
\multicolumn{5}{@{}l@{}}{\itshape 全局部件}\\
--- & 参考文献                            & ---    & ---        & 成文\\
--- & 索引                                & ---    & ---        & 成文\\

\end{longtable}
\normalsize

\vspace{0.5em}
\noindent
\textbf{关于本阶段的编号。}第五部分之前的章节尚未成文，但\emph{章号已经锁定}：
本文档当前版本只印出第 10 章，其章号即为 10 而非 1。这样做的目的是让章号成为
稳定契约 —— 后补的第 6--9 章不会使第 10 章的交叉引用失效。待 P1--P3 铺开后，
目录自会填满，而 \crefchap{ch:sp1d} 的编号不会再变。

\cleardoublepage
